\documentclass[10pt,a4paper]{amsart}
\usepackage[margin=19mm]{geometry}
\usepackage{amsmath,amssymb,mathtools}
\usepackage[scr=rsfs,cal=euler]{mathalfa}
\usepackage{xfrac}
\usepackage[unicode,hidelinks,bookmarks=false]{hyperref}
\newcommand{\ie}{i.e.\ }
\newcommand{\cf}{cf.\ }
\newcommand{\resp}{respectively\ }
\newcommand{\Subsection}{\subsection}
\providecommand{\nobreakdash}{\nobreak\hbox{-}}
\setlength{\emergencystretch}{2em}
\multlinegap0pt
\usepackage{amssymb}
\usepackage{mathtools}
\usepackage[shortlabels]{enumitem}
\usepackage{float}
\usepackage{algorithm}\usepackage{algpseudocode}
\usepackage{xcolor}
\usepackage{tikz}
\newtheorem{theorem}{Theorem}
\newcommand{\E}{\mathbb E}
\newcommand{\R}{\mathbb R}
\newcommand{\cP}{\mathcal P}
\newcommand{\cx}{\preceq_{\mathrm{cx}}}
\newcommand{\norm}[1]{\left\lVert #1\right\rVert}
\title{Barycentric Weak Inner-Product Gromov--Wasserstein}
\author{Youssef Mroueh}
\date{Source: arXiv:2608.25145v1 (CC BY 4.0)}
\begin{document}
\maketitle

\noindent\emph{Source and licence.} Barycentric Weak Inner-Product Gromov--Wasserstein. Version arXiv:2608.25145v1 (CC BY 4.0), distributed by arXiv under the Creative Commons Attribution 4.0 International licence, http://creativecommons.org/licenses/by/4.0/, as recorded in the arXiv licence metadata for this exact version and shown on the abstract page https://arxiv.org/abs/2608.25145v1. Full licence notice: \texttt{licenses/arxiv-2608.25145-CC-BY-4.0.txt}.\par\smallskip

\noindent\textit{Editorial scope.} Counted unit: the article's theorem thm:strassen (source0.tex lines 890--903). The article's complete original proof of it is delivered with it (source0.tex lines 3842--3952, one \texttt{proof} environment of 111 lines, which the article places away from the statement). No mathematics is written, repaired or re-derived: the statement and the proof are the article's own lines, in the article's own notation, and no retained line is substituted or re-flowed.  Retained uncounted as the source-local context the proof needs: lines 5242--5242.  Nothing was omitted from any retained span, so the package carries no omission record.  The retained lines cite 4 works, whose entries are transcribed verbatim from the article's own bibliography source in the e-print and delivered with short editorial labels recorded in the item's reference mapping.  No retained line contains a figure environment, an image-inclusion command or drawing code, so no image is delivered and the package declares no asset dependency.  Editorial formatting only, in addition: an AMS \texttt{amsart} preamble that restates the article's own declarations and macros used by the retained lines, namely the environments \texttt{theorem} on separate counters, together with the standard packages those lines need.  The retained environments print in this document as Theorem 1 in order of appearance, whereas the article prints the same items with the numbering of its own full document.  The complete native source of the article as distributed in the arXiv e-print for this exact version is bundled unchanged as \texttt{sources/208550/source0.tex}.
\section{Reference theorems used}\label{app:reference-theorems}

\begin{theorem}[Strassen martingale characterization]\label{thm:strassen}
For $\eta,\nu\in\cP_1(\R^d)$, $\eta\cx\nu$ if and only if there is a coupling
$\kappa\in\Pi(\eta,\nu)$ with the following martingale property. Let $(Z,Y)$ be
the coordinate random variables under $\kappa$. Then $Z\sim\eta$, $Y\sim\nu$, and
$\E_\kappa[Y\mid Z]=Z$. Equivalently, there is a
Borel probability kernel $\kappa(dy\mid z)$ satisfying
\[
  \int\kappa(\mathord\cdot\mid z)d\eta(z)=\nu,
  \qquad \int y\,\kappa(dy\mid z)=z\quad\text{for $\eta$-a.e. }z.
\]
This is Strassen's theorem \cite{Strassen}.
Appendix~\ref{app:reference-theorems} gives the complete statement, hypotheses,
and proof.
\end{theorem}

\begin{proof}
We first identify the feasible barycentric maps.  By the disintegration theorem
\cite[Theorem~3.4]{Kallenberg}, if
$\pi(dx,dy)=\mu(dx)\pi_x(dy)$ and $m_\pi(x)=\int y\,d\pi_x(y)$, then for every
integrable convex $u$, conditional Jensen gives
\[
 \int u(m_\pi(x))d\mu(x)
 \le\int\!\int u(y)d\pi_x(y)d\mu(x)=\int u(y)d\nu(y).
\]
Hence $(m_\pi)_\#\mu\cx\nu$.  Conversely, let $m_\#\mu\cx\nu$.
Theorem~\ref{thm:strassen} supplies a martingale kernel $K(z,dy)$ from
$m_\#\mu$ to $\nu$.  Define
\[
 \pi(dx,dy)=\mu(dx)K(m(x),dy).
\]
For bounded measurable $g$,
\[
 \int g(y)d\pi=\int\!\int g(y)K(z,dy)d(m_\#\mu)(z)=\int g(y)d\nu(y),
\]
and the first marginal is $\mu$ by construction.  The martingale identity gives
$\int yK(m(x),dy)=m(x)$ for $\mu$-almost every $x$.  Thus $m_\pi=m$, proving
the coupling/map equality.

The noncompact weak OT theorem
\cite[Theorems~2.9 and~3.1]{BBP} applies to the cost considered below.  It states
that if
$\nu\in\cP_2(\R^{d_y})$ and
$\widetilde C:\R^{d_x}\times\cP_2(\R^{d_y})\to\R\cup\{+\infty\}$ is jointly
lower semicontinuous for the Euclidean--$W_2$ product topology, bounded below, and
convex in its probability argument, then the weak primal admits a minimizer and
its value equals
\[
 \sup_{\psi\in\Phi_{b,2}}
 \left\{\int R_{\widetilde C}\psi\,d\mu-\int\psi\,d\nu\right\},
 \qquad
 R_{\widetilde C}\psi(x)
 :=\inf_{p\in\cP_2}\left\{\widetilde C(x,p)+\int\psi\,dp\right\},
\]
where $\Phi_{b,2}$ consists of continuous potentials bounded below by a constant
and above by a quadratic function.

For the present cost, set
\[
 C(x,p)=c(x,b(p)),\qquad b(p)=\int y\,dp(y).
\]
The barycenter map is continuous for $W_2$: if $p_n\to p$ in $W_2$, then an
optimal coupling and Cauchy--Schwarz give
$\norm{b(p_n)-b(p)}\le W_2(p_n,p)$.  Consequently $C$ is jointly continuous.
Because $b$ is affine and $c(x,\cdot)$ is convex, $p\mapsto C(x,p)$ is convex.
Let $C_0$ be the constant in the lower growth bound and put
$h(x)=C_0(1+\norm x^2)$.  Then
\[
 (C+h)(x,p)\ge\alpha\norm{b(p)}^2\ge0.
\]
Thus $C+h$ is jointly continuous, globally bounded below, and convex in $p$, while
$\nu\in\cP_2$ and $h\in L^1(\mu)$.  The cited weak OT theorem applies to
$C+h$.  Since $R_{C+h}\psi(x)=h(x)+R_C\psi(x)$, subtracting $\int h\,d\mu$ from
the primal and dual formulas yields existence of a primal minimizer and the probability-valued
dual formula for $C$.

It remains to reduce that dual to convex functions of the barycenter.  Given a
target potential $\psi$ in the weak dual, define
\[
 u(z)=\inf\left\{\int\psi(y)dp(y):p\in\cP_2(\R^{d_y}),\ b(p)=z\right\}.
\]
The competitor $p=\delta_z$ gives $u(z)\le\psi(z)$.  If $p_0,p_1$ have
barycenters $z_0,z_1$, then $(1-t)p_0+tp_1$ has barycenter
$(1-t)z_0+tz_1$; taking near minimizers proves convexity of $u$.  The growth
conditions in the weak dual make $u$ finite and lower bounded.  A finite convex
function on Euclidean space is continuous
\cite[Corollary~10.1.1]{Rockafellar}, and the Dirac competitor gives its
quadratic upper growth.  Partitioning the transform by $z=b(p)$ gives
\begin{align*}
 R_C\psi(x)
 &=\inf_{p}\left\{c(x,b(p))+\int\psi\,dp\right\}\\
 &=\inf_z\{c(x,z)+u(z)\}=Q_cu(x).
\end{align*}
For the replacement potential $u$, conditional Jensen and the Dirac competitor
give, for every $z$,
\[
 \inf_{p:\,b(p)=z}\int u(y)\,dp(y)=u(z).
\]
Consequently $R_Cu=Q_cu=R_C\psi$.  Since $u\le\psi$, the target term satisfies
$-\int u\,d\nu\ge-\int\psi\,d\nu$.  Replacing $\psi$ by $u$ therefore keeps the
transform fixed and can only increase the dual objective.  Thus convex barycentric
potentials yield the full dual value.

The dual may be indexed by all $u\in\mathcal U_2^{\rm cvx}$.  Every such $u$ has
an affine supporting lower bound
\cite[Corollary~12.1.2]{Rockafellar}.  Combining it with the coercive
quadratic lower bound on $c$ makes $Q_cu(x)>-\infty$, while evaluation at $z=0$
and the upper growth bounds give $|Q_cu(x)|\le C'(1+\norm x^2)$.  Moreover, for
every feasible $m$,
\[
 Q_cu(x)\le c(x,m(x))+u(m(x)),\qquad
 \int u(m)d\mu\le\int u\,d\nu.
\]
These inequalities give weak duality throughout $\mathcal U_2^{\rm cvx}$.  The
subclass constructed above already yields the primal value, so the supremum over
$\mathcal U_2^{\rm cvx}$ has the same value.

The feasible map set is convex: pointwise convex combinations remain below $\nu$
in convex order by Jensen.  If $c(x,\cdot)$ is $\lambda$-strongly convex and
$m_0,m_1$ are distinct minimizers, their midpoint is feasible and
\[
 \int c\left(x,\frac{m_0+m_1}{2}\right)d\mu
 \le\frac12\int c(x,m_0)d\mu+\frac12\int c(x,m_1)d\mu
 -\frac{\lambda}{8}\norm{m_0-m_1}_{L^2}^2,
\]
contradicting minimality.  Hence the optimal barycentric map is unique.
\end{proof}

\begin{thebibliography}{99}
\bibitem[BBP]{BBP} J. Backhoff-Veraguas, M. Beiglb\"ock, and G. Pammer. Existence, duality, and cyclical monotonicity for weak transport costs. \emph{Calculus of Variations and Partial Differential Equations}, 58:203, 2019.
\bibitem[Kallen]{Kallenberg} O. Kallenberg. \emph{Foundations of Modern Probability}. Third edition, Probability Theory and Stochastic Modelling 99, Springer, Cham, 2021.
\bibitem[Rockaf]{Rockafellar} R. T. Rockafellar. \emph{Convex Analysis}. Princeton Mathematical Series 28, Princeton University Press, Princeton, 1970.
\bibitem[Strass]{Strassen} V. Strassen. The existence of probability measures with given marginals. \emph{Annals of Mathematical Statistics}, 36:423--439, 1965.
\end{thebibliography}
\end{document}
